\documentclass{article}
\usepackage{graphicx} % Required for inserting images
\usepackage[margin=1.2in]{geometry} % Adjust margins
\usepackage{amsmath}
\usepackage{amssymb}
\usepackage{float}

\title{Waves}
\author{Reuben Okotie}
\date{}

\setlength{\parindent}{0pt}
\setlength{\parskip}{0.35em}

\begin{document}

\maketitle
I hope for this document to be extensive and include all the subtopics for the A-Level Physics spec. I will vet over this to have accurate notes throughout the course as a whole and hope to use this for the actual A-Level (I might also do this for select maths topics that are quite conceptually difficult like \textit{Vectors}, \textit{Differential Equations} \& \textit{Mechanics}.

For revision, attempt these Isaac Physics questions: 
\begin{itemize}
    \item https://isaacscience.org/questions/don't\_bragg\_sym
    \item Understand more about x-ray crystallography here: https://chem.libretexts.org/Bookshelves/Analytical\_Chemistry/Supplemental\_Modules\_(Analytical\_Chemistry)/Instrumentation\_and\_Analysis/Diffraction\_Scattering\_Techniques/X-ray\_Crystallography
\end{itemize}


\section{Progressive Waves}
Progressive waves are waves that transfer energy without transferring material.

\subsection{Longitudinal \& Transverse Waves}
The oscillations of transverse waves are perpendicular to the direction of energy transfer. 

The oscillations of longitudinal waves are parallel to the direction of energy transfer. 

\subsection{Polarisation}
Oscillations of transverse waves can occur at any point that is perpendicular to the direction of energy transfer; polarising the wave, thus, would ensure that the displacements are only restricted to one of these lines. 

Only transverse waves can be polarised.

Examples of polarisation: 
\begin{itemize}
    \item Polarised glasses have a polarising filter that polarises light to only one plane. This means that the glare from the surface of the water is predominately filtered out which means that the bottom of the lake can be seen with such glasses. 

    \item For old TV antennas, the polarisation of the antennas meant that only certain frequencies of the radio waves could be received and meant that there was a decreased likelihood of interference between signals - even if they were of the same frequency, allowing for targeted communications.
\end{itemize}

\section{Stationary Waves}
Stationary waves are waves that oscillate with respect to time but that do not transfer energy; their peak amplitude profiles (\textit{i.e. where the positions along the waves where peaks and troughs occur}) are fixed.

They are formed by two progressive waves moving in opposite directions interfering both constructively and destructively with the same frequency ($f$), wavelength ($\lambda$) and constant phase difference($\phi$). 

Nodes are points of total destructive interference (on stationary waves) and are points of zero displacement. Anti-nodes are points of constructive interference and are at maximum displacement (or amplitudes of the stationary wave). Anti-nodes oscillate between these points of maximum displacement.

\subsection{Harmonics}
% what is the point of harmonics (i.e. what is their importance)
Imagine a string of length L with fixed ends (nodes) at each side. The harmonics are such that integers of half a wavelength form this sequence (the harmonic sequence). The first one, where $L = \frac{\lambda}{2}$, is called the fundamental frequency, $f_0$.

The general formula for this is: 
$$L = \frac{n\lambda}{2}, \; n \in \mathbb{N}$$

When working in questions however, you may be asked to find out the frequency of a given harmonic (e.g. 5th) but can only work through it implicitly. Below is my attempt for a useful derivation of this:

It is known that $c = f\lambda$ and $L = \frac{n\lambda}{2}$. 

Therefore:
$$f = \frac{cn}{2L}$$

But given this cool derivation for $c^2 = \frac{T}{\mu}$:
https://youtu.be/1Sh7IO1EBcs

Imagine a semi-circular pulse propagating along a string and you focus on the peak of that pulse along a distance of $\Delta s$ as it propagates. 

This distance can be rearranged as, given the diagram: 
$$\Delta s = 2R\tan{\theta}$$

The resultant force at that instant (through tip-to-tail vector addition) is: 
$$\sum{F} = 2T\sin{\theta}, \; \text{where $T$ is the tension in the string}$$

Since this force is also centripetal, the resulting equation can be made:
$$2T\sin{\theta} = \frac{mc^2}{R}$$

$\mu$ is the linear density of a material and equals the mass ($m$) divided by its length ($\Delta s$). Therefore, $m = \mu\Delta s$

Putting everything into one equation:
$$2T\sin{\theta} = \frac{2\mu R \tan{\theta} c^2}{R}$$

Using the small angle approximations of $\sin{\theta} \approx \theta$ and $\tan{\theta} \approx \theta$ and further rearranging, you get to:
$$T = \mu c^2$$

Then:
$$c = \sqrt{\frac{T}{\mu}}$$

Thus the final equation for the frequency for n-th term harmonics is:
$$f = \frac{n}{2L}\sqrt{\frac{T}{\mu}}$$

\section{Refraction, diffraction \& Interference}
The \textbf{refractive index} is the factor by which the speed of light has been reduced when passing through a material. 

\textbf{Path Difference ($\Delta x$)}: The difference in length that a wave has taken to reach a given point 

\textbf{Phase Difference ($\Delta\phi$)}: The difference in phase between two given points (usually at their incidence).

Example: if wave A is at $6.56\lambda$ and wave B is at $12.31\lambda$ then:
$$\Delta x = 5.75\lambda$$
$$\Delta\phi = (.56 - .31) \times 2\pi = \frac{\pi}{2}$$

This would mean that the waves are $\frac{\pi}{2}$ out of phase. This idea is important for the upcoming experiments.

\textbf{Coherence}: When two or more waves have the same wavelength, frequency and constant phase relationships. To ensure that the waves are coherent for these upcoming experiments, monochromatic (mostly laser) light was used so that clear interference patterns could be observed. 

\subsection{Young's Double Slit Experiment}
% I will need to also include diagrams from this (probably best if drawn myself on Onenote)
Thomas Young disproved the corpuscle (\textit{i.e. light was a particle}) theory of light with this experiment.

% the experiment setup
A coherent source of light is incident on two slits which are of a distance $s$ apart, which codes for the slit separation. When the light, of wavelength $\lambda$ diffracts at these slits, they will interfere and produce a fringe pattern of alternating maxima and minima at a fringe spacing $w$ apart. The distance between the slit plane and the plane on which the interference pattern is produced is at a distance $D$ apart. 

% this needs further justification, why is this the case?
$D >> s$ such that an approximation for the path difference can be ascertained (for the derivation). 

% hand drawn diagram of this (from Onenote)
\begin{figure}[H]
    \centering
    \includegraphics[width=0.65\linewidth]{image.png}
    \caption{Diagram of the setup for this, showing the triangles needed for the later derivations (and relationships between the parameters}
    \label{fig:placeholder}
\end{figure}

As shown in the smaller green triangle:
$$\sin{\theta_1} = \frac{\lambda}{s}$$ 

The reason why that resultant length is $\lambda$ is because ... % why?

In the larger triangle: 
$$\tan{\theta_2} = \frac{w}{D}$$ % this makes more sense but I don't get why it equals $w$

Since we are using small angle approximations, both $\sin{\theta} \approx \theta$ and $\tan{\theta} \approx \theta$ and therefore we can equate both fractions such that:
$$\frac{\lambda}{s} = \frac{w}{D}$$

The desired form (in the equation sheet) is: 
$$w= \frac{w \lambda}{D}$$
% why are we using small angle approximations? I guess it's because $D >> s$
% just as a note, but for actual distance between each slit (the distance in which light can pass) what are the requiremnets there. I assume it is dependent on the wavelength for proper scattering effects to occur

\textbf{Note}: If the light was not monochromatic (e.g. a light bulb) then a filter would be made: first by shining the light through a single slit (\textbf{which achieves what effect??}) and polarising the resultant light to ensure that it is monochromatic (\textbf{again, is this right?}).

Questions are typically asked on using this experiment and transposing this knowledge with sound for example. 

\subsection{Diffraction}
\textbf{Definition}: Diffraction is the spreading out of waves. 

When there is a single slit, the diffraction pattern has the highest intensity in the middle and decreases further away from the $n=0$ order, as shown below:
\begin{figure}[H]
    \centering
    \includegraphics[width=0.65\linewidth]{image-01.png}
    \caption{Single Slit Diffraction Pattern}
    \label{fig:placeholder}
\end{figure}

The reason for diffraction when there is a single slit (and no diffraction grating) is because of Huygen's principle (\textit{... this is not in the spec ...}) where the wavefront at a boundary breaks down into an infinite set of discrete wavelet points that interfere with each other destructively and constructively to form the interference pattern as shown below: 
\begin{figure}
    \centering
    \includegraphics[width=0.95\linewidth]{image-02.png}
    \caption{Image of the intereference pattern as a result of a single slit.}
    \label{fig:placeholder}
\end{figure}

\textit{IsaacPhysics has these resources that are clearer with the theory behind Huygen's principle:}
\begin{itemize}
    \item https://isaacscience.org/concepts/cp\_reflection\_and\_refraction
    \item https://isaacscience.org/concepts/cp\_general\_waves
\end{itemize}

One issue with this experiment is that the fringe spacings are blurry (\textit{i.e. have low resolution}) and thus difficult to experimentally measure the distance. 

\textbf{Diffraction Gratings} contain a much larger amounts of minor slits (gratings) between the major slit that are closely spaced apart in order to increase the resolution of the resulting interference pattern. % why??

% derive the equation dsin\theta = n\lambda
The relationship between the order of the maxima, the angle this is away from the slit, the number of slits and the slit separation is:
$$d\sin{\theta} = n \lambda$$

% uses and examples: x-ray crystallography, spectrospocy and line spectra
\textbf{Uses of diffraction gratings}:
\begin{itemize}
    \item \textbf{X-ray crystallography}: 
\end{itemize}

\subsection{Refraction at plane surface}
\textit{I am quite confident with refraction (given the CompPhys) and, thus, will focus on total internal reflection and the nature of optics cable (and cladding and considerations there).}

\textbf{Why does the speed of light slow down in a material of higher refractive indices -- what does that actually mean?}

Total Internal Reflection occurs in a material where the incident angle is greater than the critical angle of the material. This causes the incident light (rather than refract through the material) to reflect of the normal of the material.

% Have a better appreciation for what the refractive index means and why this is the case (and kind of the theory behind this). 

This quality of materials is useful for fibre optics where signals transmitted through electromagnetic (EM) waves needs to propagate without losing data (through leakage, compression or loss of accurate encoding). Therefore, the material of fibre optics needs to be optimised to reduce these effects. 

% the type of optical fibre i need to know is a step-index optical fibre
% what are its componenents (cladding, core) and why are they optimal (and what shoudl be changed of each (relative refractive indices, sizes) to get the best out of them))

The extent of these effects are:
\begin{itemize}
    \item \textbf{Material Dispersion}: where different wavelengths of light within a material reflect at different angles since there is a slightly different associated refractive index when incident on the surface of the material. \textit{To reduce this, \textbf{using monochromatic light} a constant phase relationship of the waves as well as having them refract in the same way (since they have the same wavelength (and thus the same incident angle, rather than small deviances ($\delta\theta$) which causes the dispersion).}
    
    \item \textbf{Modal Dispersion}: Modal dispersion is the dispersion of light due to the material itself: \textit{making the diameter of the core smaller will mean that small differences in wavelength will not accumulate along the path of the fibre optic (as well as making the wire straight.}
    
    \item \textbf{Pulse Broadening}: Is the general broadening of a wave, caused by modal and material dispersion. This can cause the destructive inference of signals passing along fibre optics and thus the output signal will have a lower amplitude and longer (and inconsistent) wavelength. 
    
    \item \textbf{Absorption}: When the fibre surrounding a material absorbs its signal (parts of its energy) and thus causing a reduced amplitude emergent signal. 

% I feel there needs to be an appreciation for the differences in refractive indices for the core and cladding of a fibre optic cable. 
    
\end{itemize}
\end{document}
